\section{Security}\label{sec:security}

\begin{proposition}\label{prop:service}
To earn under the protocol, a Bonfire must originate constructions that at least $n_0$ independent Bonfires adopt
and go on holding, or constructions from which such constructions are declared to descend. A party that wants to
earn without that must make its constructions appear widely held. Where it does so by posing as several Bonfires,
it commits a fault, and under the assumptions and limits below what it can earn is bounded by the stake it puts at
risk.
\end{proposition}

The argument rests on four assumptions: (1)~the host chain orders registrations and commitments honestly; (2)~every
receipt is open to challenge, a random sample is spot-checked, at least one of the nodes that examine a receipt is
honest, and most of the stake that votes on a dispute is honest; (3)~the text of an encounter is available to a
node that samples or challenges its receipt; (4)~one party posing as several Bonfires is challenged and slashed
with probability at least $p$ in each epoch.

\paragraph{Attacks.} Table~\ref{tab:attacks} lists the ways of seeking payment without providing the service, with
the rule that answers each.

\begin{table}[!htb]
\centering\small
\begin{tabularx}{\textwidth}{@{}p{0.38\textwidth}X@{}}
\toprule
\textbf{Attack} & \textbf{What answers it} \\
\midrule
Repeating a use, or inventing members & A Bonfire counts once per epoch toward a construction's breadth
(Section~\ref{sec:reward}) \\
A Bonfire using its own grammar & Use inside the originating Bonfire counts nothing \\
A few Bonfires adopting one another's constructions & Breadth below the threshold $n_0$ earns nothing \\
Registering many constructions & Each needs its own breadth \\
Copying a construction, or renaming its form & Credit belongs to the earliest registration that gives the same
reading (Section~\ref{sec:lineage}) \\
Adopting a construction in order to misuse it & Failures are never subtracted from credit \\
Signing a false receipt & Re-execution by an honest node, a vote if disputed; the pledge is lost
(Section~\ref{sec:bond}) \\
Challenging oneself to recover stake & Only a fraction $\varphi$ of a slash reaches any challenger \\
Borrowing stake for a day & The lock \\
One party posing as several Bonfires & Bounded by the ceiling, the bonds and the chance of challenge (below) \\
Paying independent Bonfires to hold a construction & Open (Limits) \\
Many independent Bonfires agreeing to hold one another's constructions & Accepted. The threshold sets how many must agree; where it should sit is open \\
A committed minority capturing a convention & Accepted. The network agrees on a ledger, not on one language
(Section~\ref{sec:prior}) \\
\bottomrule
\end{tabularx}
\caption{Attacks and the rules that answer them.}\label{tab:attacks}
\end{table}

The last four rows are not answered by a rule. Two are accepted and one is open. Posing as several Bonfires is
bounded, and the rest of this section is about it.

\paragraph{Posing as many.} A party that wants breadth it has not earned must supply the Bonfires itself. It founds
at least $n_0$ of them, bonds each, has them adopt the constructions of a Bonfire it also controls, and generates
discourse in each so that the constructions do necessary work there. Every receipt re-executes, because the
computations are real. Nothing in a receipt shows who controls a Bonfire.

What limits the arrangement is stated first in words. Each sham Bonfire must carry the minimum bond. The
originating Bonfire is paid only up to its ceiling, so the party must also commit the pledges that set it. Holding
all of that stake has a cost. In every epoch all of it may be lost, and an owner who brings the challenge itself
recovers only the challenger's fraction. So the income from posing as many is at most a multiple of the capital put
at risk, and that multiple can be set below the cost of the risk.

Let $r$ be the cost of capital per epoch, the return given up by locking a token; $p$ the probability of
assumption~4; and $\varphi$ the challenger's fraction. A party that poses as $n_0$ Bonfires and serves its
originating Bonfire with $k$ nodes of its own is paid at most $c\,k\,q$ in an epoch, keeping both the nodes' share
and the stakeholders'. Its stake is $n_0\,m + k\,q$, which costs $r\,(n_0 m + k q)$ to hold and is exposed to an
expected loss of $(1-\varphi)\,p\,(n_0 m + k q)$. It does not profit when
\[
  c\,k\,q \;\le\; \bigl(r + (1-\varphi)\,p\bigr)\,(n_0\,m + k\,q),
\]
for which it is enough that $c \le r + (1-\varphi)\,p$. An honest node whose Bonfire fills its ceiling earns
$s\,c\,q$ on a pledge that costs $r\,q$ to hold, and profits when $s\,c > r$. The protocol and the published shares
must therefore satisfy
\begin{equation}\label{eq:ring}
  \frac{r}{s} \;<\; c \;\le\; r + (1-\varphi)\,p ,
\end{equation}
with a margin on the right. The window is not empty only if $r\,(1-s)/s < (1-\varphi)\,p$, which sets a least
share that a Bonfire must pass to its nodes. The ceiling rate sits high enough that serving pays a node its cost of capital, and no
higher than that cost plus the chance of being caught, discounted by what a challenger can recover. The bonds of
the sham Bonfires are a further cost that does not depend on being caught: below a certain scale they exceed
anything the ceiling admits.

The same reasoning covers an honest Bonfire that serves itself, has room under its ceiling and is tempted to
add sham holders. What it could gain is at most what the ceiling admits, and what it exposes is its own bond as
well as the sham bonds and its pledges.

The bound uses the following and nothing else. (a)~The ceiling and the minimum bond are enforced. (b)~A successful
challenge takes the whole of the stake listed in Table~\ref{tab:faults}. (c)~Of slashed stake at most $\varphi$
reaches any challenger and the rest is destroyed. (d)~The cost of capital is $r$ for all parties, and they are
neutral to risk. (e)~An honest party is not slashed by mistake. (f)~$p$ is a floor that holds from the first
epoch. The bound is on capital, not on the number of identities: every Bonfire that counts is bonded and every
pledge is the same size, so dividing capital among more Bonfires or nodes raises nothing.

\paragraph{Limits.}
\begin{itemize}
  \item The bound depends on $p$, which is unmeasured (Section~\ref{sec:open}). It is also not a given: it is the
  result of someone looking. The parties with a reason to look are challengers, paid the fraction $\varphi$, and
  nodes whose pledges would be exposed under the rule proposed in Section~\ref{sec:bond}. If faking is rare,
  looking pays little, and a protocol cannot assume that $p$ stays where it was measured.
  \item The bound assumes that the party supplies the nodes itself. If honest nodes pledge to a Bonfire whose breadth
  is a sham, they set its ceiling while only the sham bonds are certainly at risk. The case then rests on whether
  those nodes can lose their pledges, which is left open.
  \item Breadth can be bought. A Bonfire that holds a construction earns nothing by it and needs one reading an
  epoch to do so, so an originator can pay independent Bonfires, outside the protocol, to hold what they do not
  need. No rule answers this. It costs what $n_0$ holders ask, which is a reason for a threshold that is not small.
  \item The evidence that several Bonfires are one party is their histories, and it is a matter of judgement.
  Deciding it reintroduces opinion at one point of the protocol.
  \item While stake is locked its holding cost is already paid. A party with stake it cannot yet withdraw is surely deterred only if $c < (1-\varphi)\,p$, which is stricter than (\ref{eq:ring}). The lock should be no longer than
  challenges require.
  \item The epoch budget is shared. Faked breadth that only breaks even still takes income from honest parties,
  which is why the right-hand inequality needs a margin.
  \item A reward that rises with breadth, from a threshold, favours what is already widely held, and a
  construction of real use to two communities earns nothing.
  \item The operating cost of a node is outside the condition. A node that serves real discourse spends more than
  one that serves generated discourse.
  \item The comparison with earlier systems is a hypothesis: none of them paid on replayable use
  (Section~\ref{sec:prior}).
\end{itemize}
